\section{Recovery Quality and Robustness}
\label{sec:families}
\sectionaim{Test whether recovery is safe, tail-aware, and stable across builds.}

\subsection{RQ5: Can Target Recalibration Recover Value?}
\textbf{Question.} Can a small amount of target-side information repair unsafe source reuse without returning to the fixed endpoint? \textbf{Protocol.} The target-selection lane shifts the frozen query-indexed action pool over the registered grid, selects on 500 queries, certifies on a disjoint 500, and evaluates on 1,000. \textbf{Observation.} Risk falls from \WZeroSiftSourceRisk\% to \WZeroSiftTargetRisk\% on SIFT and from \WZeroArxivSourceRisk\% to \WZeroArxivTargetRisk\% on Arxiv. Mean gains remain \WZeroSiftTargetGain\% and \WZeroArxivTargetGain\%, with registered intervals [\WZeroSiftTargetGainLow, \WZeroSiftTargetGainHigh]\% and [\WZeroArxivTargetGainLow, \WZeroArxivTargetGainHigh]\%. \textbf{Conclusion.} Target evidence restores both empirical safety and substantial search savings on both datasets; the result supports recalibration for recurring profiled queries, not free transfer to unseen queries.

\subsection{RQ6: Does Recovery Damage the Tail?}
\textbf{Question.} Are mean gains purchased by worse high-cost queries or pervasive fallback? \textbf{Protocol.} We compare query-pooled p95 and p99 distance counts against the same fixed endpoint and retain fallback and nonfallback outcomes. \textbf{Observation.} Target-recalibrated p95 is \WZeroSiftTargetPninetyfive\ versus \WZeroSiftTargetFixedPninetyfive\ on SIFT and \WZeroArxivTargetPninetyfive\ versus \WZeroArxivTargetFixedPninetyfive\ on Arxiv. Registered p99 comparisons have the same non-worsening direction. Fallback occurs on 0/10 SIFT and 1/10 Arxiv targets, so the aggregate tail result is not produced by replacing most targets with the endpoint. \textbf{Conclusion.} Recovery improves the observed p95 and does not introduce a structural p99 or fallback-tail regression in this replay. These are distance-count comparisons, not wall-clock noninferiority certificates.

\begin{table}[t]
\caption{Tail and fallback evidence for target recalibration.}
\label{tab:tail}
\small
\begin{tabular*}{\linewidth}{@{\extracolsep{\fill}}lrrr@{}}
\toprule
Dataset & Target p95 & Endpoint p95 & Fallback \\
\midrule
SIFT & \WZeroSiftTargetPninetyfive & \WZeroSiftTargetFixedPninetyfive & \WZeroSiftTargetFallbackBuilds/10 \\
Arxiv & \WZeroArxivTargetPninetyfive & \WZeroArxivTargetFixedPninetyfive & \WZeroArxivTargetFallbackBuilds/10 \\
\bottomrule
\end{tabular*}
\end{table}

\subsection{RQ7: Is the Result Driven by One Build?}
\textbf{Question.} Could one favorable target account for the recovery gain? \textbf{Protocol.} We combine the registered build-outer nested bootstrap with LOTO and deletion of the largest-gain build. \textbf{Observation.} The gain interval remains strictly positive on both datasets. Minimum LOTO and delete-largest gains are 44.37\% on SIFT and 44.23\% on Arxiv, closely tracking the pooled 44.39\% and 44.79\%. \textbf{Conclusion.} No single target build explains the mean benefit. The claim is build-robust within the ten registered targets; it does not replace validation on independently sampled deployment populations.
